% collectionview-compositional.tex — UICollectionViewCompositionalLayout: item → group →
% section → layout, dimensions, orthogonal scrolling, supplementary/decoration items,
% section providers, list configuration, cell registration + content configuration,
% self-sizing, invalidation. Diffable data sources live on tableview-diffable.tex.
% Sources (checked 2026-10-05):
%   developer.apple.com/documentation/uikit/uicollectionviewcompositionallayout (iOS 13)
%   developer.apple.com/documentation/uikit/nscollectionlayoutgroup (repeatingSubitem:count:, iOS 16)
%   developer.apple.com/documentation/uikit/uicollectionlayoutsectionorthogonalscrollingbehavior
%   developer.apple.com/documentation/uikit/uicollectionlayoutsectionorthogonalscrollingproperties (iOS 17)
%   developer.apple.com/documentation/uikit/uicollectionlayoutlistconfiguration (iOS 14)
%   developer.apple.com/documentation/uikit/uicollectionview/cellregistration (iOS 14) ·
%   .../uicollectionviewcell/configurationupdatehandler (iOS 15) · .../uihostingconfiguration (iOS 16)
%   developer.apple.com/documentation/uikit/uicollectionviewdiffabledatasource/sectionidentifier(for:) (iOS 15)
%   WWDC21 10252 "Make blazing fast lists and collection views" (registration inside the provider throws)
% @labels: area=ios-platform kind=api platform=ios topic=ui
% @tags: uicollectionview, compositional-layout, nscollectionlayoutgroup, fractional-dimensions, orthogonal-scrolling, section-provider, list-configuration, cellregistration, content-configuration, supplementary-views, self-sizing
\documentclass[8pt]{extarticle}
\usepackage{printup-sheet}
\usepackage{array}

\lstdefinelanguage{SwiftSheet}{
  morekeywords={protocol,class,final,struct,enum,func,var,let,weak,init,
    if,else,return,guard,self,nil,try,await,async,throws,private,some,static,
    switch,case,in,true,false,AnyObject,Void,String,Bool},
  sensitive=true, morecomment=[l]{//}, morecomment=[s]{/*}{*/}, morestring=[b]"}

\tikzset{
  lbl/.style={font=\tiny, text=black!75, inner sep=1pt, align=center},
  lay/.style={draw=sheetGrey, thick, rounded corners=2pt, fill=black!2},
  sec/.style={draw=sheetBlue, thick, rounded corners=2pt, fill=sheetBlue!6},
  grp/.style={draw=sheetGreen, thick, rounded corners=1.5pt, fill=sheetGreen!8},
  itm/.style={draw=sheetOrange, thick, rounded corners=1pt, fill=sheetOrange!12},
  sup/.style={draw=sheetBrown, dashed, fill=sheetBrown!10},
  card/.style={draw=sheetBlue, fill=sheetBlue!15, rounded corners=1pt},
}

\begin{document}

\sheettitle{UICollectionView — compositional layout, registrations, configurations}{ios-platform · folio}

\oneliner{A compositional layout (iOS 13) is \textbf{declared}, not computed in delegate
callbacks: \textbf{items} sit in \textbf{groups} (the repeating unit), a group repeats
along a \textbf{section}, sections stack into the \textbf{layout}. Every size is
\texttt{.absolute}, \texttt{.estimated} or \texttt{.fractional} \emph{of its container}.
Cells are dequeued from \textbf{registrations} and rendered from \textbf{content
configurations} (iOS 14).}

\vspace{3pt}
\noindent\begin{tikzpicture}[sheet]
  % ── left: the nesting ──
  \draw[lay] (0,-0.35) rectangle (7.5,3.3);
  \node[lbl, anchor=north west] at (0.05,3.28) {\textbf{layout} — sections stack along the scroll axis};
  \draw[sec] (0.2,-0.2) rectangle (7.3,2.95);
  \node[lbl, anchor=north west] at (0.25,2.93) {\textbf{section} — \texttt{contentInsets} · \texttt{interGroupSpacing} · \texttt{orthogonal…}};
  \draw[sup] (0.4,2.2) rectangle (7.1,2.55);
  \node[lbl] at (3.75,2.375) {\textbf{header} = boundary supplementary, \texttt{.top}, \texttt{pinToVisibleBounds}};
  % group
  \draw[grp] (0.4,0.3) rectangle (5.15,2.0);
  \node[lbl, anchor=north west] at (0.45,1.98) {\textbf{horizontal group} w \texttt{1.0} · h \texttt{.absolute(120)}};
  \foreach \i in {0,1,2} {
    \draw[itm] ({0.55+\i*1.55},0.45) rectangle ({1.9+\i*1.55},1.6);
    \node[lbl] at ({1.225+\i*1.55},1.15) {\textbf{item}};
    \node[lbl] at ({1.225+\i*1.55},0.8) {w \texttt{1/3} of\\the \emph{group}};}
  % badge on item 3
  \draw[sheetRed, fill=sheetRed!60] (5.03,1.6) circle (0.09);
  \node[lbl, anchor=west, text=sheetRed] at (5.2,1.62) {badge};
  % next group
  \draw[grp, dashed] (5.35,0.3) rectangle (7.1,1.35);
  \node[lbl] at (6.22,0.82) {next group $\to$\\repeats};
  \node[lbl, text=sheetBrown, anchor=west] at (0.25,0.05) {section fill = \textbf{decoration} (no data, registered on the layout)};
  % ── middle: the screen ──
  \draw[thick, rounded corners=4pt, sheetGrey] (7.9,-0.35) rectangle (10.9,3.3);
  % section 0 carousel
  \node[lbl, anchor=west] at (7.95,3.1) {\textbf{s0} carousel};
  \draw[card] (8.0,2.2) rectangle (8.2,2.95);
  \draw[card, fill=sheetBlue!30] (8.3,2.2) rectangle (10.5,2.95);
  \draw[card] (10.6,2.2) rectangle (10.8,2.95);
  \draw[<->, thick, sheetOrange] (8.4,2.08) -- (10.4,2.08);
  \node[lbl, text=sheetOrange] at (9.4,1.93) {orthogonal paging};
  % section 1 grid
  \node[lbl, anchor=west] at (7.95,1.72) {\textbf{s1} 2-col grid};
  \foreach \r in {0,1} \foreach \c in {0,1}
    \draw[itm] ({8.05+\c*1.4},{0.95-\r*0.6}) rectangle ({9.35+\c*1.4},{1.5-\r*0.6});
  % section 2 list
  \node[lbl, anchor=west] at (7.95,0.1) {\textbf{s2} list};
  \foreach \r in {0,1} \draw[sheetGrey, fill=white] (8.05,{-0.1-\r*0.1}) rectangle (10.75,{-0.02-\r*0.1});
  \draw[->, thick, sheetGrey] (11.05,3.1) -- node[lbl, right, align=left]{main\\scroll} (11.05,-0.2);
  % ── right: dimensions ──
  \node[font=\bfseries\small, anchor=west] at (11.75,3.15) {\texttt{NSCollectionLayoutDimension}};
  \node[lbl, anchor=north west, align=left, text width=48mm] at (11.75,2.85) {
    \texttt{.absolute(120)} — fixed points\\[1pt]
    \texttt{.estimated(44)} — a guess; the cell's Auto Layout measures the real size\\[1pt]
    \texttt{.fractionalWidth(0.5)} / \texttt{Height} — of the \textbf{container}: an
    item's container is its \textbf{group}; a group's is the section's content area
    (after insets). Height may be \texttt{.fractionalWidth} — squares.\\[3pt]
    \textbf{orthogonal:} \texttt{.none} · \texttt{.continuous} ·
    \texttt{.continuousGroupLeadingBoundary} · \texttt{.paging} (container width) ·
    \texttt{.groupPaging} · \texttt{.groupPagingCentered} (one group)\\[3pt]
    \textbf{s0} = group width \texttt{0.9} so neighbours peek; at \texttt{1.0} it looks like
    \texttt{.paging}.};
\end{tikzpicture}

\begin{multicols}{2}
\raggedright

\section{How it works}
\begin{itemize}
  \item \textbf{Group} — \texttt{.horizontal}/\texttt{.vertical(layoutSize:subitems:)},
        \texttt{repeatingSubitem:count:} (iOS 16), groups nest (big tile + 2×2);
        \texttt{.custom} returns explicit frames. \texttt{interItemSpacing}
        (\texttt{.fixed}/\texttt{.flexible}) between items only.
  \item \textbf{Section provider} — the closure init \texttt{\{ index, env in … \}}
        builds each section; \texttt{env.container} (effective content size) +
        \texttt{env.traitCollection} pick 2 vs 4 columns. Map index $\to$ your section
        via \texttt{sectionIdentifier(for:)} (iOS 15), not raw indices.
  \item \textbf{List} (iOS 14) — \texttt{UICollectionLayoutListConfiguration}:
        \texttt{.plain}/\texttt{.grouped}/\texttt{.insetGrouped}/\texttt{.sidebar},
        \texttt{headerMode}, swipe-action providers;
        \texttt{.list(using:layoutEnvironment:)} per section mixes lists with grids —
        the modern \texttt{UITableView}.
  \item \textbf{Supplementary} — headers/footers = boundary items (\texttt{elementKind},
        \texttt{alignment}, \texttt{pinToVisibleBounds}); badges = item-anchored
        (\texttt{containerAnchor}); data-backed, via \texttt{SupplementaryRegistration}.
        \textbf{Decoration} = visual only, class registered on the layout.
  \item \textbf{Registration} (iOS 14) — \texttt{CellRegistration<Cell, Item>} +
        \texttt{dequeueConfiguredReusableCell(using:for:item:)}: no string ids, typed
        item. Create it \textbf{once}.
  \item \textbf{Content configuration} — \texttt{defaultContentConfiguration()} $\to$
        \texttt{UIListContentConfiguration} (\texttt{text}, \texttt{secondaryText},
        \texttt{image}) $\to$ \texttt{cell.contentConfiguration}: a value describes the
        look, the content view applies it. Custom: \texttt{UIContentConfiguration}
        (\texttt{makeContentView()}, \texttt{updated(for:)}) + a \texttt{UIContentView}.
        States: \texttt{configurationUpdateHandler} (iOS 15). SwiftUI:
        \texttt{UIHostingConfiguration} (iOS 16).
  \item \textbf{Self-sizing} — \texttt{.estimated} on \textbf{both} item and group for
        that axis + a fully constrained cell; estimates near reality or content jumps.
  \item \textbf{Invalidation} — a bounds change re-runs the provider (rotation for
        free); \texttt{invalidateLayout()} for other inputs;
        \texttt{setCollectionViewLayout(\_:animated:)} swaps layouts.
\end{itemize}

\section{Flow layout vs compositional}
{\footnotesize
\begin{tabular}{@{}>{\raggedright\arraybackslash}p{20mm}>{\raggedright\arraybackslash}p{25mm}>{\raggedright\arraybackslash}p{28mm}@{}}
\toprule
 & flow layout & compositional \\ \midrule
sizes & delegate \texttt{sizeForItemAt} & declared dimensions \\
a section scrolls sideways & no & \texttt{orthogonal…} \\
table-style lists & no — \texttt{UITableView} & \texttt{.list(using:)} \\
badges, backgrounds & subclass the layout & supplementary · decoration \\
adapt to width & invalidate, recompute & provider gets \texttt{env} \\
\bottomrule
\end{tabular}}

\columnbreak

\section{Example}
\begin{lstlisting}[language=SwiftSheet]
let layout = UICollectionViewCompositionalLayout { i, env in
  if i > 0 { return .list(using: .init(appearance: .insetGrouped),
                          layoutEnvironment: env) } // a list
  let item = NSCollectionLayoutItem(layoutSize: .init(
    widthDimension: .fractionalWidth(1),
    heightDimension: .fractionalHeight(1)))  // fill the group
  let group = NSCollectionLayoutGroup.horizontal(layoutSize:
    .init(widthDimension: .fractionalWidth(0.9),   // peek
          heightDimension: .absolute(180)), subitems: [item])
  let s = NSCollectionLayoutSection(group: group)
  s.orthogonalScrollingBehavior = .groupPagingCentered
  return s }
let reg = UICollectionView.CellRegistration<      // ONCE
  UICollectionViewListCell, Item> { cell, _, item in
  var c = cell.defaultContentConfiguration()
  c.text = item.title; cell.contentConfiguration = c }
\end{lstlisting}

\section{Traps}
\begin{itemize}
  \trap{\texttt{sizeForItemAt} does nothing — only the flow layout reads
        \texttt{UICollectionViewDelegateFlowLayout}. The silent migration bug.}
  \trap{\texttt{.fractionalHeight(1)} on an item = fill the \textbf{group}, not the screen.}
  \trap{Item \texttt{contentInsets} (shrink each item, gaps double between neighbours) +
        \texttt{interItemSpacing} = surprise gutters. Pick one; section insets for edges.}
  \trap{\texttt{.estimated} on the item but \texttt{.absolute} on the group — no
        self-sizing.}
  \trap{Creating a \texttt{CellRegistration} inside the cell provider — re-registers per
        cell; iOS 15+ raises an exception.}
  \trap{An orthogonal section scrolls in its own internal scroll view —
        \texttt{scrollViewDidScroll} does not fire; use
        \texttt{visibleItemsInvalidationHandler} (per frame: keep it cheap).}
\end{itemize}

\section{Timeline}
\begin{tikzpicture}[sheet]
  \tikzset{yr/.style={font=\bfseries\scriptsize, text=sheetBlue},
           ev/.style={font=\tiny, align=center, text=black!85, text width=15mm, anchor=north}}
  \draw[thick, sheetGrey, ->] (0,0) -- (8.3,0);
  \foreach \x/\y/\t in {
      0.8/iOS 13/{compositional layout\\diffable data source},
      2.45/iOS 14/{list configuration\\registrations\\content configs},
      4.1/iOS 15/{\texttt{configuration-}\\\texttt{UpdateHandler}\\\texttt{reconfigureItems}},
      5.75/iOS 16/{\texttt{UIHosting-}\\\texttt{Configuration}\\\texttt{repeatingSubitem}},
      7.4/iOS 17/{orthogonal\\scrolling\\properties}}{
    \fill[sheetBlue] (\x,0) circle (0.6mm);
    \node[yr, above=1pt] at (\x,0) {\y};
    \node[ev] at (\x,-0.12) {\t};
  }
\end{tikzpicture}

\section{Remember}
\textbf{Item in group in section in layout; fractions of the container; estimate both
levels; register once, configure by value.}

\end{multicols}

\noindent{\footnotesize\color{sheetGrey}\textit{Related:} table view + diffable data
source (snapshots, \texttt{reconfigureItems}) · auto layout (self-sizing) · rendering
pipeline · SwiftUI \texttt{List}/\texttt{LazyVGrid} · accessibility}

\end{document}
